\documentclass[11pt,a4paper]{article}
\usepackage[bahasa]{babel}
\usepackage[utf8]{inputenc}
\usepackage[T1]{fontenc}
\usepackage{geometry}
\geometry{a4paper, margin=2.2cm, top=2.5cm, bottom=2.5cm}
\usepackage{xcolor}
\usepackage[most]{tcolorbox}
\usepackage{booktabs}
\usepackage{tabularx}
\usepackage{enumitem}
\usepackage{amsmath}
\usepackage{amssymb}
\usepackage{hyperref}
\usepackage{fancyhdr}
\usepackage[expansion=false]{microtype}

% Warna PLN
\definecolor{plnblue}{RGB}{0,75,135}
\definecolor{plnorange}{RGB}{255,107,53}
\definecolor{plngreen}{RGB}{0,150,136}

\hypersetup{
    colorlinks=true,
    linkcolor=plnblue,
    urlcolor=plnblue,
    citecolor=plnblue
}

% Header and footer
\pagestyle{fancy}
\fancyhf{}
\fancyhead[L]{\small\textbf{Handout Diklat Modul 4.3}}
\fancyhead[R]{\small PT PLN (Persero)}
\fancyfoot[C]{\small\thepage}
\renewcommand{\headrulewidth}{0.4pt}

% Box styles
\newtcolorbox{plnbox}[1][]{%
  colback=plnblue!5!white,
  colframe=plnblue,
  fonttitle=\bfseries,
  title=#1,
  sharp corners,
  boxrule=0.8pt,
  breakable
}

\newtcolorbox{warnbox}[1][]{%
  colback=plnorange!10,
  colframe=plnorange,
  fonttitle=\bfseries,
  title=#1,
  sharp corners,
  boxrule=0.8pt,
  breakable
}

\newtcolorbox{outputbox}[1][]{%
  colback=plngreen!10,
  colframe=plngreen!70!black,
  fonttitle=\bfseries,
  title=#1,
  sharp corners,
  boxrule=0.8pt,
  breakable
}

\setlist[itemize]{leftmargin=*, topsep=2pt, itemsep=2pt}
\setlist[enumerate]{leftmargin=*, topsep=2pt, itemsep=2pt}

\begin{document}

% ================= COVER =================
\begin{titlepage}
\centering
\vspace*{1.2cm}
{\Large\bfseries\color{plnblue} PROGRAM PENGUATAN KOMPETENSI}\\[0.2cm]
{\Large\bfseries\color{plnblue} TATA KELOLA KEARSIPAN}\\[0.8cm]

{\Huge\bfseries\color{plnblue} HANDOUT DIKLAT}\\[0.3cm]
{\LARGE\bfseries Modul 4.3}\\[0.3cm]
{\Large Optimalisasi Prosedur dan Alur Kerja Terintegrasi}\\[0.8cm]
\rule{0.82\linewidth}{0.6pt}\\[0.6cm]

{\large Level 4 (Mastery)}\\
{\large Target Peserta: Manajemen Atas (MA)}\\[0.4cm]
{\large PT PLN (Persero)}

\vfill
\begin{plnbox}[Tujuan Handout]
Handout ini merangkum materi inti Bab 1 sampai Bab 6 sebagai panduan cepat saat pembelajaran, diskusi, dan tindak lanjut implementasi di unit kerja.
\end{plnbox}

\vfill
{\small \today}
\end{titlepage}

\tableofcontents
\newpage

% ================= BAGIAN AWAL =================
\section{Ringkasan Modul}
\begin{plnbox}[Fokus Strategis]
Modul ini menempatkan kearsipan sebagai \textit{value center}, bukan sekadar \textit{cost center}. Tujuan akhirnya adalah proses kearsipan yang cepat, patuh regulasi, teraudit, dan siap mendukung keputusan bisnis PLN.
\end{plnbox}

\subsection{Kompetensi Akhir Peserta}
Peserta ditargetkan mampu:
\begin{enumerate}
    \item Mengevaluasi proses \textit{as-is} dan menemukan area optimasi.
    \item Merancang workflow \textit{to-be} berbasis integrasi digital.
    \item Menyusun SOP digital terintegrasi yang sah dan operasional.
    \item Menghitung efisiensi dan ROI sebagai justifikasi bisnis.
\end{enumerate}

\subsection{Output Portofolio Wajib}
\begin{outputbox}[Empat Dokumen Hasil Diklat]
\begin{enumerate}
    \item Matriks Identifikasi Area Optimasi.
    \item Diagram BPMN 2.0 To-Be + peta pemicu integrasi.
    \item Draft SOP Digital Terintegrasi.
    \item KPI Tracker + Action Plan 30 Hari.
\end{enumerate}
\end{outputbox}

\subsection{Alokasi Waktu Pembelajaran}
\begin{table}[h!]
\centering
\small
\begin{tabularx}{\textwidth}{@{} l X c @{} }
\toprule
\textbf{Bagian} & \textbf{Aktivitas} & \textbf{Durasi} \\
\midrule
Fondasi dan Diagnosa & Pengantar, teori, evaluasi proses, pemetaan masalah & 25 menit \\
Desain dan Implementasi & Rancang To-Be, SOP digital, KPI-ROI, workshop & 50 menit \\
\midrule
\textbf{Total} & \textbf{Materi dan praktik} & \textbf{80 menit} \\
\bottomrule
\end{tabularx}
\end{table}

\newpage

% ================= BAB 1 =================
\section{Bab 1 - Pendahuluan}
\subsection{Inti Materi}
\begin{itemize}
    \item Kearsipan PLN bersifat strategis untuk keandalan operasional, kepatuhan, dan tata kelola.
    \item Tantangan utama saat ini: proses manual, duplikasi input, metadata tidak konsisten, dan bottleneck persetujuan.
    \item Arah perbaikan: prosedur terintegrasi yang mendukung \textit{single source of truth}.
\end{itemize}

\subsection{Pesan Kunci Manajerial}
\begin{plnbox}[Kenapa Ini Urgen]
Setiap keterlambatan dan ketidakakuratan arsip berpotensi menjadi risiko bisnis, risiko audit, dan risiko hukum. Karena itu, optimalisasi prosedur adalah agenda kepemimpinan, bukan sekadar agenda administrasi.
\end{plnbox}

% ================= BAB 2 =================
\section{Bab 2 - Evaluasi Prosedur dan Identifikasi Area Optimasi}
\subsection{Inti Materi}
\begin{itemize}
    \item Evaluasi dilakukan dengan basis fakta lapangan, bukan asumsi.
    \item Teknik utama: \textit{document shadowing}, \textit{log analysis}, dan \textit{gemba walk}.
    \item Hasil evaluasi harus dituangkan ke Matriks Identifikasi Area Optimasi.
\end{itemize}

\subsection{Kerangka Diagnosa 7 Domain}
\begin{table}[h!]
\centering
\small
\begin{tabularx}{\textwidth}{@{} c X X @{} }
\toprule
\textbf{No} & \textbf{Domain} & \textbf{Pertanyaan Kendali} \\
\midrule
1 & Kebijakan & Apakah kebijakan internal selaras dengan regulasi? \\
2 & Tata Kelola & Apakah peran dan akuntabilitas jelas? \\
3 & SDM & Apakah kompetensi digital arsip memadai? \\
4 & Interoperabilitas & Apakah data mengalir lintas sistem tanpa input ulang? \\
5 & Proses Bisnis & Apakah ada langkah non-value added? \\
6 & Keamanan & Apakah akses, jejak audit, dan integritas terjaga? \\
7 & Preservasi & Apakah retensi, disposisi, dan keberlanjutan arsip terkelola? \\
\bottomrule
\end{tabularx}
\end{table}

% ================= BAB 3 =================
\section{Bab 3 - Perancangan Alur Kerja Terintegrasi (To-Be)}
\subsection{Prinsip Desain}
\begin{enumerate}
    \item \textbf{Single Entry, Multiple Use}: data dimasukkan sekali, dipakai lintas proses.
    \item \textbf{Rule-Based Routing}: jalur persetujuan dikunci aturan sistem.
    \item \textbf{Explicit Exception Handling}: ada mekanisme fallback dan eskalasi.
\end{enumerate}

\subsection{BPMN untuk Manajemen}
Gunakan BPMN 2.0 untuk memvalidasi alur lintas fungsi. Fokus simbol manajerial:
\begin{itemize}
    \item \textbf{Start Event}: pemicu proses.
    \item \textbf{Task}: pekerjaan sistem/manusia.
    \item \textbf{Gateway}: keputusan dan percabangan.
    \item \textbf{Sequence/Message Flow}: urutan dan komunikasi lintas sistem.
    \item \textbf{End Event}: titik akhir proses dan status arsip.
\end{itemize}

\begin{warnbox}[Batasan]
Bab ini tidak membahas detail teknis API, database, atau middleware. Fokus pada logika bisnis, otorisasi, dan titik kontrol arsip.
\end{warnbox}

% ================= BAB 4 =================
\section{Bab 4 - Penyusunan SOP Digital Terintegrasi}
\subsection{Inti Materi}
SOP digital harus mengikat tiga unsur: proses bisnis, kontrol kepatuhan, dan kontrol sistem.

\subsection{Komponen Minimal SOP}
\begin{enumerate}
    \item Tujuan dan ruang lingkup.
    \item Definisi dan istilah kerja.
    \item Peran, kewenangan, dan otorisasi.
    \item Alur proses (rujukan BPMN To-Be).
    \item Standar metadata dan validasi.
    \item Ketentuan TTE tersertifikasi.
    \item Ketentuan audit trail dan log aktivitas.
    \item SLA, eskalasi, dan penanganan pengecualian.
\end{enumerate}

\subsection{Checklist Kepatuhan Dokumen Digital}
\begin{itemize}
    \item Dokumen memiliki identitas unik, waktu, dan sumber penciptaan.
    \item Dokumen disahkan melalui mekanisme TTE tersertifikasi.
    \item Riwayat akses/perubahan terekam dan dapat ditelusuri.
    \item Integritas dokumen dapat diverifikasi (misal hash).
    \item Retensi dan disposisi mengacu jadwal retensi arsip.
\end{itemize}

% ================= BAB 5 =================
\section{Bab 5 - Penghitungan Efisiensi dan ROI}
\subsection{KPI Utama}
\begin{table}[h!]
\centering
\small
\begin{tabularx}{\textwidth}{@{} X X @{} }
\toprule
\textbf{Indikator} & \textbf{Target Pasca Optimasi} \\
\midrule
Waktu Siklus Proses & turun 40--60\% \\
First-Time Accuracy Metadata & minimal 90\% \\
Waktu Temu Kembali Arsip Digital & kurang dari 30 detik \\
Kepatuhan SLA dan Jejak Audit & minimal 95\% \\
\bottomrule
\end{tabularx}
\end{table}

\subsection{Rumus Dasar ROI Waktu}
\[
\text{ROI Waktu} = (\text{Waktu Lama} - \text{Waktu Baru}) \times \text{Volume Tahunan} \times \text{Biaya per Jam/Hari}
\]

\subsection{Kerangka Nilai Time-Cost-Risk}
\begin{itemize}
    \item \textbf{Time Saving}: pemangkasan durasi proses.
    \item \textbf{Cost Saving}: penurunan biaya cetak, kirim, simpan, rework.
    \item \textbf{Risk Saving}: penghindaran risiko audit, sanksi, dan sengketa.
\end{itemize}

\begin{plnbox}[Catatan Implementasi]
Gunakan data riil unit kerja saat menyusun business case. Angka simulasi di modul berfungsi sebagai contoh pembelajaran, bukan angka audit resmi.
\end{plnbox}

% ================= BAB 6 =================
\section{Bab 6 - Penutup dan Refleksi}
\subsection{Ringkasan Capaian}
Pada akhir modul, peserta diharapkan siap memimpin pilot implementasi optimalisasi prosedur kearsipan di unitnya masing-masing.

\subsection{Rencana Tindak Lanjut 30 Hari}
\begin{enumerate}
    \item Minggu 1: validasi proses \textit{as-is} dan tetapkan prioritas area optimasi.
    \item Minggu 2: finalisasi BPMN To-Be dan aturan routing persetujuan.
    \item Minggu 3: selesaikan draft SOP digital + checklist kepatuhan.
    \item Minggu 4: presentasi business case (KPI dan ROI) untuk persetujuan pilot.
\end{enumerate}

\begin{outputbox}[Komitmen Eksekusi]
Keberhasilan modul diukur dari implementasi nyata di unit kerja: proses lebih cepat, kepatuhan lebih kuat, dan bukti arsip lebih dapat dipertanggungjawabkan.
\end{outputbox}

\newpage
\section{Lampiran Ringkas}
\subsection{Template Cepat Matriks Area Optimasi}
\begin{table}[h!]
\centering
\small
\begin{tabularx}{\textwidth}{@{} p{0.7cm} X X X p{2.2cm} @{} }
\toprule
\textbf{No} & \textbf{Temuan As-Is} & \textbf{Dampak} & \textbf{Akar Masalah} & \textbf{Prioritas} \\
\midrule
1 & & & & \\
2 & & & & \\
3 & & & & \\
4 & & & & \\
\bottomrule
\end{tabularx}
\end{table}

\subsection{Template Cepat KPI Tracker}
\begin{table}[h!]
\centering
\small
\begin{tabularx}{\textwidth}{@{} X c c c @{} }
\toprule
\textbf{Indikator} & \textbf{Baseline} & \textbf{Target} & \textbf{Realisasi} \\
\midrule
Waktu Siklus & & & \\
First-Time Accuracy & & & \\
Waktu Temu Kembali & & & \\
Kepatuhan SLA/Trail & & & \\
\bottomrule
\end{tabularx}
\end{table}

\vspace{0.4cm}
\begin{center}
\small\textit{Selesai. Gunakan handout ini bersama modul bab, slide presentasi, dan narasi fasilitator untuk implementasi terarah.}
\end{center}

\end{document}
